% Fichier complété par tools/complete_missing_corrections.py.
% Source : STAT/STAT-02-Local/1023_Vilebrequin/1023_Vilebrequin.tex
% Statut : rédigé (4 question(s)).
\begin{corrigebox}[Corrigé rédigé]
\CorrigeQuestion{1}
On décompose la pièce en volumes élémentaires. Le centre d'inertie est
                $\overrightarrow{OG}=\sum m_i\overrightarrow{OG_i}/\sum m_i$; les trous
                sont traités comme des masses négatives.
\CorrigeQuestion{2}
Pour que $G$ appartienne à l'axe $(O,\vec x)$, ses coordonnées
                transversales doivent être nulles. L'équation barycentrique en $y$ (et en
                $z$ si nécessaire) fournit directement la hauteur $h$.
\CorrigeQuestion{3}
On remplace les masses par $\rho V_i$; la masse volumique commune se
                simplifie. L'application numérique est alors effectuée avec les dimensions
                du dessin.
\CorrigeQuestion{4}
En $G$, $\{T_{p\to S}\}=\{-mg\vec y_0;\vec0\}_G$. En $O$ :
                $\{T_{p\to S}\}=\{-mg\vec y_0;
                \overrightarrow{OG}\wedge(-mg\vec y_0)\}_O$.
\end{corrigebox}
